% Gabarit pandoc des corrigés publiés (brut et guidé). Contenu sans marque ni nom de site.
% Compilation : lualatex (latexmk -lualatex).
% Le Markdown reste la source unique : ne pas éditer ce fichier .tex à la main.
\documentclass[11pt,a4paper]{article}

% --- Moteur, polices (Latin Modern) et langue ---------------------------------
\usepackage{iftex}
\ifLuaTeX\else\errmessage{Ce document se compile avec lualatex}\fi
\usepackage{amsmath,amssymb,mathtools}
\usepackage{fontspec}
% mathrm=sym : \mathrm prend ses glyphes dans la police mathématique (ℓ, μ, Ω), comme KaTeX.
\usepackage[math-style=TeX,bold-style=TeX,mathrm=sym]{unicode-math}
% Repli pour les glyphes absents de Latin Modern (indices et exposants Unicode,
% flèches, filets des schémas en texte, symboles) : DejaVu, puis Latin Modern Math
% (ℓ, lettres mathématiques Unicode des légendes).
\directlua{
  luaotfload.add_fallback("repliTexte", {"DejaVuSerif:mode=harf;", "DejaVuSans:mode=harf;", "file:latinmodern-math.otf:mode=harf;"})
  luaotfload.add_fallback("repliMono", {"DejaVuSansMono:mode=harf;", "DejaVuSans:mode=harf;", "file:latinmodern-math.otf:mode=harf;"})
}
\setmainfont{Latin Modern Roman}[RawFeature={fallback=repliTexte}]
\setsansfont{Latin Modern Sans}[RawFeature={fallback=repliTexte}]
\setmonofont{Latin Modern Mono}[RawFeature={fallback=repliMono}]
\setmathfont{Latin Modern Math}
\usepackage[french]{babel}
\frenchsetup{SmallCapsFigTabCaptions=false}

% --- Mise en page --------------------------------------------------------------
\usepackage[a4paper,margin=2.2cm,headheight=14pt,footskip=1.2cm]{geometry}
\usepackage{microtype}
\setlength{\emergencystretch}{4em}
\tolerance=1500
\setlength{\parindent}{0pt}
\setlength{\parskip}{0.55\baselineskip plus 2pt minus 1pt}
\usepackage{xcolor}
% Gris nommés : babel-french rend « ! » actif dans le document.
\definecolor{gris2}{gray}{0.98}
\definecolor{gris3}{gray}{0.97}
\definecolor{gris5}{gray}{0.95}
\definecolor{gris35}{gray}{0.65}
\definecolor{gris40}{gray}{0.60}
\definecolor{gris45}{gray}{0.55}
\definecolor{gris50}{gray}{0.50}
\definecolor{gris55}{gray}{0.45}
\definecolor{gris60}{gray}{0.40}
\definecolor{gris65}{gray}{0.35}
\definecolor{gris70}{gray}{0.30}
\definecolor{gris80}{gray}{0.20}
\definecolor{gris85}{gray}{0.15}
\usepackage{graphicx}
\usepackage{float}
\usepackage{caption}
\captionsetup{font=small,labelformat=empty,justification=centering}
\usepackage{booktabs,longtable,array,multirow,calc}
\usepackage{siunitx}
\sisetup{locale=FR}
\usepackage[version=4]{mhchem}
\usepackage{listings}
\usepackage{fancyhdr}
\usepackage{lastpage}
\usepackage[most]{tcolorbox}
\usepackage{enumitem}
\setlist{nosep,leftmargin=*,topsep=2pt}
\usepackage[hidelinks,unicode]{hyperref}
\usepackage{bookmark}
\hypersetup{pdftitle={Centrale-Supélec TSI — Rédaction 2026},pdfauthor={Rédigé par IA},pdfcreator={LaTeX},pdfsubject={Corrigé brut}}
\usepackage{xurl}

% --- Pandoc : listes, tableaux, figures ----------------------------------------
\providecommand{\passthrough}[1]{#1}
\providecommand{\tightlist}{\setlength{\itemsep}{0pt}\setlength{\parskip}{0pt}}

\makeatletter
\def\maxwidth{\ifdim\Gin@nat@width>\linewidth\linewidth\else\Gin@nat@width\fi}
\def\maxheight{\ifdim\Gin@nat@height>0.55\textheight 0.55\textheight\else\Gin@nat@height\fi}
\makeatother
\setkeys{Gin}{width=\maxwidth,height=\maxheight,keepaspectratio}
\setcounter{secnumdepth}{-\maxdimen}
\providecommand{\pandocbounded}[1]{#1}

% --- Code (listings) -----------------------------------------------------------
\lstset{
  basicstyle=\ttfamily\small,
  columns=fullflexible,
  keepspaces=true,
  breaklines=true,
  breakatwhitespace=false,
  postbreak=\mbox{\textcolor{gris50}{$\hookrightarrow$}\space},
  frame=single,
  rulecolor=\color{gris35},
  framesep=4pt,
  xleftmargin=5pt,xrightmargin=5pt,
  backgroundcolor=\color{gris3},
  keywordstyle=\bfseries,
  commentstyle=\color{gris65},
  stringstyle=\color{gris80},
  showstringspaces=false,
  tabsize=4,
  extendedchars=true,
  aboveskip=6pt,belowskip=6pt,
}

% --- Macros des corrigés (KaTeX ↔ LaTeX) ---------------------------------------
% Les corrigés sont écrits pour KaTeX : on définit ici ce que KaTeX connaît et
% que LaTeX n'a pas, sans jamais retoucher le Markdown.
\providecommand{\R}{\mathbb{R}}
\providecommand{\N}{\mathbb{N}}
\providecommand{\Z}{\mathbb{Z}}
\providecommand{\Q}{\mathbb{Q}}
\providecommand{\C}{\mathbb{C}}
\providecommand{\dd}{\mathrm{d}}
\providecommand{\llbracket}{\lBrack}
\providecommand{\rrbracket}{\rBrack}
\providecommand{\lBrace}{\lbrace}
\providecommand{\rBrace}{\rbrace}
\providecommand{\Box}{\square}
\providecommand{\Diamond}{\lozenge}
\providecommand{\varnothing}{\emptyset}
\providecommand{\argmax}{\operatorname*{arg\,max}}
\providecommand{\argmin}{\operatorname*{arg\,min}}
\renewcommand{\leq}{\leqslant}
\renewcommand{\geq}{\geqslant}

% Équation hors texte : réduite à la largeur de la ligne si elle déborde.
\newsavebox{\boiteequation}
\newcommand{\equationajustee}[1]{%
  \sbox{\boiteequation}{$\displaystyle #1$}%
  \ifdim\wd\boiteequation>\linewidth
    \[\resizebox{\linewidth}{!}{\usebox{\boiteequation}}\]%
  \else
    \[#1\]%
  \fi}

% --- Encadrés pédagogiques (sobres, lisibles en noir et blanc) -----------------
\tcbset{
  encadre/.style={
    enhanced,breakable,sharp corners,
    colback=white,colframe=gris55,
    boxrule=0.4pt,leftrule=2.2pt,
    left=6pt,right=6pt,top=3pt,bottom=3pt,
    before skip=6pt,after skip=6pt,
    fonttitle=\bfseries\small,coltitle=black,
    attach title to upper,after title={\quad},
    parbox=false,
  }
}
\newtcolorbox{indication}[1]{encadre,colframe=gris45,colback=gris2,title={Indication~#1}}
\newtcolorbox{methode}{encadre,colframe=gris70,title={Méthode}}
\newtcolorbox{piege}{encadre,colframe=gris85,boxrule=0.8pt,leftrule=2.2pt,title={Piège}}
\newtcolorbox{jury}{encadre,colframe=gris60,colback=gris5,title={Jury}}
\newtcolorbox{encadrenote}{encadre,colframe=gris40}

% --- Titres de section ---------------------------------------------------------
\usepackage{titlesec}
\titleformat{\section}{\Large\bfseries}{}{0pt}{}
\titleformat{\subsection}{\large\bfseries}{}{0pt}{}
\titleformat{\subsubsection}{\normalsize\bfseries}{}{0pt}{}
\titlespacing*{\section}{0pt}{14pt plus 4pt minus 2pt}{6pt}
\titlespacing*{\subsection}{0pt}{10pt plus 3pt minus 2pt}{4pt}
\titlespacing*{\subsubsection}{0pt}{8pt plus 2pt minus 2pt}{2pt}

% --- En-tête et pied de page ---------------------------------------------------
\pagestyle{fancy}
\fancyhf{}
\fancyhead[L]{\small\textsc{Corrigé brut} · Centrale-Supélec TSI}
\fancyhead[R]{\small Rédaction 2026}
\fancyfoot[L]{\footnotesize Corrigé rédigé par IA, non officiel}
\fancyfoot[R]{\footnotesize \thepage\,/\,\pageref*{LastPage}}
\renewcommand{\headrulewidth}{0.4pt}
\renewcommand{\footrulewidth}{0pt}
\fancypagestyle{plain}{\fancyhf{}\fancyfoot[R]{\footnotesize \thepage\,/\,\pageref*{LastPage}}\fancyfoot[L]{\footnotesize Corrigé rédigé par IA, non officiel}\renewcommand{\headrulewidth}{0pt}}

\begin{document}
\thispagestyle{plain}
\begin{center}
  {\small\scshape Corrigé brut}\par\smallskip
  {\LARGE\bfseries Centrale-Supélec TSI — Rédaction 2026\par}\medskip
  {\large Centrale-Supélec · TSI · Rédaction · 2026\par}
\end{center}
\vspace{2pt}\hrule\vspace{8pt}

\emph{Corrigé rédigé par IA, non officiel. Signalez toute erreur.}

Épreuve de 4 heures, sans calculatrice ni document. Thème 2025-2026 : « Expériences de la nature » ; œuvres au programme : Jules Verne, \emph{Vingt mille lieues sous les mers} ; Georges Canguilhem, \emph{La Connaissance de la vie} (extraits au programme) ; Marlen Haushofer, \emph{Le Mur invisible}.

Convention de décompte retenue : est compté comme un mot toute unité séparée des autres par une espace, un trait d'union ou une apostrophe d'élision (« l'imprévu » = 2 mots, « c'est-à-dire » = 4 mots, « aujourd'hui » = 1 mot).

\hypertarget{partie-a-ruxe9sumuxe9-de-texte}{%
\section{Partie A --- Résumé de texte}\label{partie-a-ruxe9sumuxe9-de-texte}}

\textbf{Problématique du texte :} pourquoi le jardinage soulage-t-il le mal de vivre ? Parce que le jardin rend visibles l'interdépendance des vivants, la résilience que donne la diversité et le surgissement imprévu de la vie, qui appellent l'humilité.

Cultiver un jardin soulage, malgré ses peines physiques : le jardinage guérit du mal de vivre. Pourquoi ? Parce que le jardinier y recouvre la puissance, ou du moins foi en la vie. Chaque être œuvre et tout est lié : aucun vivant n'est vain. Si l'on tient chacun pour un \textbf{\textbar{}} sujet, selon la symbiose réciproque de Serres, un équilibre s'installe malgré la prédation. Plus un jardin est divers, plus il résiste, ses espèces se soutenant face aux perturbations. Qu'il reste le même quand ses espèces changent se discute, mais la diversité fonde la résilience, et ce qui semblait \textbf{\textbar{}} mort revit. Surgit aussi l'imprévu, souvent heureux : il suscite l'admiration cartésienne puis oblige à s'adapter. Jardiner, c'est donc agir humblement, car nul ne lutte longtemps contre la nature. Loin des idéologies, on affronte le réel : la vie jaillit, avec et sans nous, et ce surgissement console. \textbf{\textbar{}}

\textbf{Total : 150 mots.}

\hypertarget{partie-b-dissertation}{%
\section{Partie B --- Dissertation}\label{partie-b-dissertation}}

\textbf{Problématique :} l'expérience de la nature livre-t-elle un ordre déjà là, que le regard n'aurait qu'à recueillir, ou le sens et la valeur qu'on lui prête sont-ils institués par le vivant qui regarde, au prix d'une violence que la formule passe sous silence ?

Le jardin est un morceau de nature enclos, assez petit pour qu'un regard en fasse le tour, assez vivant pour qu'il s'y passe toujours quelque chose. C'est de ce lieu que Sonia Feertchak tire une évidence : on y découvre que « tout fait sens, l'harmonie se donne à voir, rien de vivant n'est insignifiant ». Trois affirmations s'y superposent. La première est intellectuelle : la nature est intelligible, chaque être y tient une place et une fonction. La deuxième est sensible : cet ordre n'a pas à être démontré, il s'offre de lui-même au regard. La troisième est morale : aucun vivant n'est négligeable. Or les trois œuvres au programme mettent précisément à l'épreuve l'expérience de la nature : un savant embarqué derrière la vitre du salon du Nautilus dans \emph{Vingt mille lieues sous les mers} de Jules Verne, un philosophe qui interroge dans \emph{La Connaissance de la vie} les conditions auxquelles le vivant peut être connu, une femme enfermée avec quelques bêtes au fond d'une vallée autrichienne dans \emph{Le Mur invisible} de Marlen Haushofer. L'expérience de la nature livre-t-elle un ordre déjà là, que le regard n'aurait qu'à recueillir, ou le sens et la valeur qu'on lui prête sont-ils institués par le vivant qui regarde, au prix d'une violence que la formule passe sous silence ? Les trois œuvres donnent d'abord raison à Feertchak ; elles récusent ensuite l'idée d'une harmonie simplement donnée ; il apparaît alors que la troisième proposition seule résiste, à condition de l'entendre comme une exigence du regard plutôt que comme un constat.

L'expérience de la nature commence bien par une révélation d'ordre. Le roman de Verne est bâti sur ce dispositif : le grand panneau vitré du salon fait de l'océan un spectacle continu, et Aronnax, professeur au Muséum, y reconnaît, nomme et classe. Conseil, son domestique, décline embranchements, classes, ordres, genres et espèces ; le « monstre » insaisissable des premiers chapitres, que son maître prenait pour un narval gigantesque, se résout lui-même en machine identifiable. La mer apparaît surtout comme un système où rien ne se perd. Nemo s'habille du byssus des pinnes marines, fume une algue riche en nicotine, ne se nourrit que de ce que l'océan produit ; au « royaume du corail », ce sont les polypes qui scellent la tombe du matelot mort, la nature achevant le travail des hommes. Tout, jusqu'au cadavre, entre dans un cycle : l'ordre est intelligible parce qu'il est intégral.

Canguilhem donne à cette intuition son fondement philosophique. Dans « Le vivant et son milieu », il rappelle, après Uexküll, qu'un vivant n'habite pas l'univers indifférent de la physique mais un milieu propre, où les excitations valent comme des signaux : pour l'animal, rien n'est neutre, tout lui est adressé. Dans « Le normal et le pathologique », il montre que la vie n'est pas indifférente à ses conditions : le vivant est normatif, il préfère certains états à d'autres et institue ses propres normes. Là où Feertchak écrit que rien de vivant n'est insignifiant, Canguilhem permet de préciser que rien n'est insignifiant pour un vivant, puisque vivre c'est déjà valoriser.

Haushofer fait de cette découverte la matière même de son récit. La narratrice du \emph{Mur invisible} est une citadine sans expérience de la campagne ; l'enfermement la contraint à voir. Elle apprend les saisons, la fenaison, la culture des pommes de terre, l'estivage à l'alpage avec la vache Bella ; elle nomme, soigne et protège le chien Lynx, la vieille chatte, le veau, et pleure les chatons perdus. Peu à peu, elle cesse de se penser au centre et se découvre une créature parmi les autres. Le jardin de Feertchak, la vallée de Haushofer et le salon du Nautilus sont ainsi trois enclos qui rendent l'attention possible, et où l'ordre du vivant commence à se laisser lire.

Mais la formule dit davantage : l'harmonie « se donne à voir », comme si l'ordre s'offrait indépendamment de l'œil qui le reçoit. C'est là que les œuvres résistent. Verne dispose devant la même vitre trois regards incompatibles : Conseil classe sans s'émouvoir, Ned Land juge les poissons à leur qualité comestible, Aronnax s'émerveille en savant. Le spectacle est un ; le sens qu'on y lit dépend de celui qui regarde. Canguilhem l'établit en droit : le milieu propre d'un animal, l'environnement géographique et l'univers de la science ne coïncident pas, et le décentrement scientifique ruine l'évidence d'une harmonie visible. Une harmonie « à hauteur de jardinier » reste une harmonie pour le jardinier. « L'expérimentation en biologie animale » ajoute une objection décisive : on ne connaît le vivant qu'en intervenant sur lui, en l'isolant, en le séparant de la totalité qu'il forme, c'est-à-dire en altérant ce que l'on prétend observer. La contemplation pure est un mythe : il n'y a pas de spectacle du vivant sans dispositif, et le hublot du Nautilus en est un.

Surtout, cet ordre a un prix, que les deux romans font payer. Nemo refuse que Ned Land harponne les baleines franches, inoffensives et inutiles à l'équipage, car ce serait tuer pour tuer ; puis il lance l'éperon du Nautilus contre une troupe de cachalots qu'il juge malfaisants, et le récit d'Aronnax devient celui d'un carnage. La hiérarchie que l'harmonie semblait dissoudre renaît sous le nom de nuisible : l'insignifiant revient toujours par là. Chez Haushofer, la narratrice doit tuer pour vivre ; elle abat des chevreuils, les dépèce, et sa répugnance ne se résorbe jamais dans l'idée d'un équilibre. Le mur lui-même constitue l'objection majeure : rien ne l'explique, aucune cause n'en est jamais donnée, et de l'autre côté un vieil homme reste figé près de sa fontaine. L'expérience de la nature peut donc confronter à l'absolument insensé. On mesure alors ce que le propos de Feertchak doit à son cadre : le jardin est un lieu protégé, entretenu, choisi ; la vallée est une prison et l'océan un abîme. L'harmonie ne se donne pas : elle se construit, et elle suppose des conditions que l'expérience brute de la nature ne garantit pas.

Reste la troisième proposition, et c'est elle que les œuvres valident, mais en la déplaçant. Canguilhem lui donne un fondement dans « La monstruosité et le monstrueux » : la monstruosité, non la mort, est la contre-valeur vitale, car elle révèle que la forme normale était une réussite et non une nécessité ; et si le fantastique est un monde, la vie, elle, est pauvre en monstres. « Machine et organisme » complète l'analyse : la machine exécute un plan qui lui vient du dehors, tandis que l'organisme, polyvalent, improvise et répare. Il n'y a donc de réussite et d'échec, de normal et de monstrueux, que pour le vivant. La formule de Feertchak ne décrit pas un fait de la nature : elle énonce une propriété du vivant comme tel, être de valeur qui ne peut être compté pour rien sans que l'on manque ce qu'il est.

Haushofer en donne l'image la plus juste avec la corneille blanche. Rejetée par les siennes, sans fonction dans aucune économie naturelle, elle ne sert à rien ; la narratrice la nourrit pourtant chaque jour. L'oiseau n'est pas significatif en soi : l'attention le rend tel. De même, la fin du récit, où un inconnu tue à coups de hache le taureau puis Lynx avant d'être abattu par la narratrice, est insoutenable parce que ces vies-là n'étaient pas interchangeables ; c'est le lien tissé, non l'ordre du monde, qui interdit de les tenir pour rien. Le récit lui-même, écrit pour personne, est l'acte par lequel elle arrache à l'oubli ce que nul ne viendra compter.

Verne opère le même déplacement. Quand Nemo affirme que la mer est l'immense désert où l'homme n'est jamais seul, car il sent frémir la vie à ses côtés, il ne constate pas une plénitude : il l'affirme contre l'apparence du désert, en homme qui a rompu avec la société. Les ruines de l'Atlantide, entrevues à la lueur d'un volcan sous-marin, montrent d'ailleurs que la nature recouvre les ordres humains sans les commenter, et le Maelström final, où le Nautilus disparaît aux yeux du narrateur, rappelle que l'ordre déborde toujours celui qui l'observe. Le dernier mot peut revenir à Canguilhem : la pensée du vivant doit tenir du vivant l'idée du vivant. Reconnaître qu'aucun vivant n'est insignifiant, ce n'est pas lire une harmonie dans les choses, c'est se savoir soi-même vivant parmi les vivants, ce que Feertchak appelle, à la fin de son texte, l'humilité.

Des trois affirmations, la première n'est donc vraie qu'en partie : la nature fait sens pour un vivant, non en soi. La deuxième ne résiste pas : l'harmonie ne se donne pas, elle se conquiert par un dispositif, une attention, et parfois une violence assumée. La troisième est vraie, mais comme exigence et non comme description. L'expérience de la nature n'est pas le spectacle d'un ordre ; elle est le travail par lequel un vivant en reconnaît d'autres et se découvre situé parmi eux. Le jardinier de Feertchak, le professeur Aronnax et la femme du \emph{Mur invisible} font, chacun à son échelle, cette même expérience, à ceci près que le jardinier peut refermer la barrière derrière lui.

\textbf{Total : environ 1600 mots} (sous le plafond de 1800 mots).

\end{document}
